\section{Further Preliminaries}
\label{sec:futher_prelim}
In this section, we present additional statistical and \gls{ml} concepts that we will encounter throughout the thesis. 


\subsection{Importance Sampling}
\label{subsec:importance_sampling}

\Glspl{is}~\cite{mcbook} is a statistical tool to estimate the expectation $\mu = \expectedvalue_{x \sim P}[f(x)]$ of some function $f:X\mapsto \realnumbers$ under a \emph{target} distribution $P$ while having access only to samples collected with another distribution $Q$ --the \emph{behavioural} distribution. 
We note $p$ and $q$ the density functions of $P$ and $Q$, respectively.
If $P$ is absolutely continuous \wrt $Q$, which we note $P \ll Q$, then importance sampling builds an unbiased estimator of $\mu$ in the following way,
\begin{align*}
    \widehat{\mu} = \frac{1}{N}\sum_{i=1}^{N} \frac{p(x_{i})}{q(x_{i})}f(x_{i})
\end{align*}
where $\{x_{i}\}_{i=1}^{N} \sim Q$. 

To go further, {\glspl{mis}} considers a set of behavioural distributions $(Q_j)_{j\in[\![1,J]\!]}$ instead of the unique distribution $Q$. 
We note $(q_j)_{j\in[\![1,J]\!]}$ the density functions of $(Q_j)_{j\in[\![1,J]\!]}$.
We note $N$ be the total number of samples and $N_j$ the ones sampled by $Q_j$, which therefore satisfy $\textstyle N=\sum\nolimits_{j=1}^{J} N_j$. 
If $P \ll Q_j$ for all $j\in[\![1,J]\!]$, then the unbiased \gls{mis} estimator reads,
\begin{align*}
	\widehat{\mu} = \sum_{j=1}^J \frac{1}{N_j}\sum_{i=1}^{N_j} \beta_j(x_{ij})\frac{p(x_{ij})}{q_j(x_{ij})}f(x_{ij}),
\end{align*}
where $\{x_{ij}\}_{i=1}^{N_j} \sim Q_j$ and $(\beta_j(x))_{j\in[\![1,J]\!]}$
are a partition of the unit for every $x \in X$.
The choice of this partition is free, but a common choice is the \emph{balance heuristic} (BH)~\cite{veach1995monte}, which sets
\begin{align*}
    \beta_j(x)=\frac{N_j q_j(x)}{\sum_{k=1}^J N_k q_k(x)}.
\end{align*}
BH can be seen as casting the estimator back to the classic importance sampling estimator, where the samples can be regarded as coming from the mixture:
\begin{align*}
    \Phi = \sum_{k=1}^J\frac{N_k}{N}Q_k.
\end{align*} 


\subsection{Divergences}
The divergence in statistics defines a notion of distance between probability distributions. 
We define two such divergences, which we use in this work.

\noindent\textbf{Kullback-Leibler divergence~\cite[Section~2.3]{cover1991information}.}\indent 
For two probability distributions $P$ and $Q$ on the same probability space and such that $P \ll Q$, the Kullback-Leibler divergence between them is defined as,
\begin{align*}
    \kullbackleiblerdivergence(P\Vert Q)=\int_{\Omega} p(\omega) \log \frac{p(\omega)}{q(\omega)} \de \omega.
\end{align*}

\noindent\textbf{\Renyi divergence~\cite{renyi1961measures}.}\indent The $\alpha$-\Renyi divergence between two probability distributions $P$ and $Q$ such that $P \ll Q$ is defined for $\alpha \in [0,\infty]$ as:
\begin{align}
\label{eq:renyi_div}
	\renyidivergence[\alpha](P\Vert Q) = \frac{1}{\alpha-1}\log \int_{\Omega} p(\omega)^\alpha q(\omega)^{1-\alpha} \de \omega.
\end{align}
We note $\exponentialrenyidivergence[\alpha](P\Vert Q) = \exp\{\renyidivergence[\alpha](P\Vert Q)\}$ the exponential $\alpha$-\Renyi divergence.
The \Renyi divergence is interesting as it is related to the $\alpha$-moment of the importance weight in the following way,
\begin{align*}
    \expectedvalue_{x\sim Q}\left[\left(\frac{p(x)}{q(x)}\right)^{\alpha}\right] = \exponentialrenyidivergence[\alpha](P\Vert Q)^{\alpha-1}.
\end{align*}
Note that the 1-\Renyi divergence, which is not defined by \Cref{eq:renyi_div} but obtained by continuity \cite{van2014renyi}, corresponds to the Kullback-Leibler divergence.
In this work, we will always consider the 2-\Renyi divergence, and for simplicity, we will drop the exponent in the notation, $\renyidivergence[2]=\renyidivergence$.



\subsection{Neural Networks}
A \gls{nn} is a parameterised function $f$ that results from the stacking of affine mappings combined with point-wise non-linear functions. 
The latter plays a central role in the ability of a neural network to approximate complex functions. They are called \emph{activation functions}. 
Common activation functions are the sigmoid or rectified linear unit (ReLU). 
Numerous architectures for neural networks have been proposed over the years, from simpler \emph{fully-connected} networks (known as well as dense, linear or feedforward networks) to more complex deep neural networks \cite{goodfellow2016deep} with more intricate structure and/or more layers and parameters.
The fully connected network consists of a matrix of parameters $W\in\realnumbers^{n\times m}$, also called \emph{weights}, and an activation function $\sigma$.
When applied to some input $x\in\realnumbers^m$, the network outputs $y=\sigma(Wx)\in\realnumbers^n$.
If only a single combination of linear and non-linear functions is applied, the network is actually called a \emph{layer}. 
Another network can be obtained by stacking several such layers. 
A second ubiquitous type of layer is \emph{convolution}. 
A convolution scans through the input--in one, two, or even more dimensions--by recursively applying a kernel to only a subset of this input. 
The size of the kernel is referred to as the \emph{receptive field}. 
Because the kernel is fixed, this operation is akin to the mathematical definition of the convolution, hence its name. 
The great advantages of convolutions over fully-connected layers are the reduced number of weights and the structured mapping of the input. 
This explains the great empirical success of convolutions \cite{krizhevsky2017imagenet}. 
When applied to \emph{time series}--sequences indexed by time--convolution can proceed in a way that preserves causality and is referred to as \emph{temporal convolutions}~\cite{oord2016wavenet, liotet2020deep}.
We now delve into two more complex networks that will be used later in this dissertation.


\subsubsection{Transformers}
\label{subsubsec:transformer}
The \emph{Transformer} \cite{vaswani2017attention} is a network that has been introduced for sequence-to-sequence tasks, that is, encoding a sequence to then decode it into another sequence. 
The Transformer uses the principle of attention \cite{bahdanau2014neural} coupled with feed-forward layers to encode or decode a sequence.
On top of its better performance, one main advantage of the Transformer compared to previous approaches for sequence-to-sequence tasks is its computational cost. 
Previous approaches relied on \emph{recurrent neural networks} that process the input sequentially and perform \emph{back-propagation through time}, which greatly slows the process.
Another useful property of Transformers is that they can be implemented in a way that preserves the causality of the input sequence.  
Before concluding this short presentation, we focus on one specific layer inside the Transformer, the \emph{positional encoding}.
It will be particularly useful later on.
Positional encoding is an early layer of the Transformer that allows one to add information on the position of an element inside the sequence.
It computes a vectorial representation of the index of this element as a set of sinusoidal transformations applied to this index. 
An interesting property of this representation is that, however large the value of the index can be, the representation's output is sinusoidal and is therefore bounded.
Moreover, the size of the vectorial representation can be controlled by the network designer.


\subsubsection{Flow-Based Models}
\label{subsubsec:maf}

Flow-based models are a type of \emph{generative model}, such as Generative Adversarial Networks (GAN) \cite{goodfellow2014generative}. Generative models are \glspl{nn} trained on a dataset $\mathcal{D}$ to generate new data points with the same statistics as $\mathcal{D}$.
However, contrarily to GANs and most generative models, flow-based ones not only learn this generative model, but also learn the original probability $p$ used to generate $\mathcal{D}$. 
This property, together with the great empirical results of flow-based approaches \cite{oord2016wavenet,kingma2018glow}, is why we will use them in this dissertation.

In particular, we will present the \gls{maf} \cite{papamakarios2017masked} network. It combines two ideas,  \emph{normalizing flow}~\cite{rezende2015variational} and \emph{autoregressive density estimation}~\cite{uria2016neural}.
The idea of autoregressive density estimation methods such as Masked Autoencoder Distribution Estimator (MADE)~\cite{germain2015made} on which MAF is based is to generate a sample recursively, dimension after dimension. 
Practically, to sample a vector $\vectorialform{x}\in\realnumbers^n$ from a probability distribution $p$, autoregressive density estimation rewrites $\textstyle p(x)=\prod\nolimits_{i=1}^{n}p(x_i\vert x_{1:i-1})$ and each conditional $p(x_i\vert x_{1:i-1})$ can be approximated by a \gls{nn}.
Normalising flows instead learn to generate a sample $\vectorialform{x}\sim p$ by applying some invertible function $f$ to a vector $\vectorialform{u}\in\realnumbers^m$ sampled from a base distribution $\rho$. 
This base distribution is usually simple, such as a normal distribution.
The probability obtained can be expressed as
\begin{align}
\label{eq:p_inverted_maf}
    p(\vectorialform{x}) = \rho(f^{-1}(\vectorialform{x})) \left\lvert \det\left(\frac{\partial f^{-1}}{\partial \vectorialform{x}}\right)\right\rvert.
\end{align}
MAF combines these approaches and generates a sample $\vectorialform{x}$ dimension-wise as,
\begin{align}
\label{eq:sampling_maf}
    x_i = u_i \exp(\alpha_i) + \mu_i 
\end{align}
where $\mu_i=f_{\mu_i}(x_{1:i-1})$ and $\alpha_i=f_{\alpha_i}(x_{1:i-1})$ are modelled by neural networks and $u_i\sim\mathcal{N}(0,1)$.

As anticipated, MAF learns the density of a data point $\vectorialform{x}$, and it is possible to access it using \Cref{eq:p_inverted_maf}. 
Due to the recursive structure of the generating process, the inverse function in this equation has a simple expression: $\textstyle\exp\left(\sum\nolimits_{i=1}^n\alpha_i\right)$\cite{papamakarios2017masked}.
Normalising flow layers can be stacked to approximate more complex probabilities; the same idea can be applied to MAF.

Calling $p_{\vectorialform{\theta}}$ the probability it has learnt, the MAF network is trained to minimise the Kullback-Leibler divergence $\kullbackleiblerdivergence(p,p_{\vectorialform{\theta}})$ which is estimated using samples from the training set $\mathcal{D}$. 
In practice, the network is thus trained on the following objective,
\begin{align*}
    \underset{\vectorialform{\theta}}{\text{maximise}}   \sum_{\vectorialform{x}\in\mathcal{D}}\log p_{\vectorialform{\theta}} (\vectorialform{x}).
\end{align*}